Additive models with pairwise interactions are widely used for interpretable machine learning, yet their component functions suffer from a fundamental identifiability problem: variance can be freely shifted between main effects and interactions without changing predictions, yielding contradictory explanations for the same model.
Lengerich et al.\ (2020) proposed a purification algorithm based on the functional analysis of variance that enforces integrate-to-zero constraints via iterative mass-moving, producing a unique, canonical decomposition.
This report reproduces their method for piecewise-constant models (depth-2 XGBoost ensembles) and extends the original work with two new experiments.
First, we apply purification to the German Credit dataset under three weighting distributions (empirical, uniform, and independence-assumed), confirming that predictions remain invariant (RMSE\,=\,2043.3 across all modes) while the attribution of variance to main effects versus interactions changes with the choice of weights.
Second, we conduct a synthetic stress-test with controlled feature dependence ($\rho \in \{0, 0.3, 0.6, 0.9, 0.97\}$) and two ground-truth functions, measuring weighted $L_2$ distance, sign flip rate, and interaction variance share against the empirical-weight baseline.
At $\rho = 0.97$, misspecified weights produce interaction surfaces with $L_2$ distances up to 0.90, sign flip rates of 17--34\%, and variance attribution shifts from near 0\% (empirical) to 73\% (independence-assumed).
These results demonstrate that while purification successfully resolves the identifiability problem, the resulting explanations are highly sensitive to the assumed weighting distribution when features are dependent, underscoring the need for careful weight specification and sensitivity analysis in practice.
